


本文围绕抱轨式磁浮列车悬浮电磁铁功率放大器故障检测问题展开研究，目标是建立电磁力预测模型、车体--悬浮架垂向动力学模型，并完成整车功率放大系数估计与单电磁铁故障定位。针对四个问题，本文依次采用\textbf{最小二乘拟合}、\textbf{二自由度垂向动力学建模}、\textbf{动力学反演}和\textbf{相对异常特征检测}的方法，实现从电磁力建模到故障诊断的完整分析流程。

针对问题一，根据电磁力与电流平方成正比、与悬浮间隙平方成反比的物理规律，并考虑电流方向对电磁力方向的影响，建立单电磁铁电磁力预测模型
$\hat F=K_0I|I|/z^2$。利用附件1中 $99999$ 组有效样本进行最小二乘拟合，得到 $K_0=0.0799874103\,\mathrm{N\cdot m^2/A^2}$，后续计算中近似取 $\hat F=0.0800I|I|/z^2$。模型决定系数为 $0.878800$，均方根误差为 $1346.6659\,\mathrm{N}$，相对均方根误差为 $4.9915\%$，残差均值接近零，说明模型具有较好的预测精度和物理可解释性。

针对问题二，本文以车体和悬浮架为研究对象，建立空气弹簧--阻尼连接下的\textbf{车体--悬浮架二自由度垂向动力学模型}，并将附件2中16台电磁铁实际电磁力之和作为外部输入进行数值积分。考虑自由动力学模型可能导致悬浮间隙超出物理范围，进一步引入完整硬约束 $0\leq z_{\mathrm{gap}}(t)\leq0.06\,\mathrm{m}$。计算结果表明，$t=9\,\mathrm{s}$ 时悬浮架位移为 $z_f(9)=0.0005295023\,\mathrm{m}$，对应悬浮间隙为 $z_{\mathrm{gap}}(9)=0.0594704977\,\mathrm{m}$，即 $59.4705\,\mathrm{mm}$。

针对问题三，在16台电磁铁功率放大系数相同且固定的假设下，先由问题一模型计算理想总电磁力，再由车体加速度和悬浮间隙二阶导数反演实际总电磁力，并采用零截距最小二乘法估计统一功率放大系数。全样本估计值为 $\hat{\eta}_{\mathrm{all}}=0.7071311922$，考虑二阶微分端点误差，剔除首尾 $0.10\,\mathrm{s}$ 后得到主要估计值 $\hat{\eta}\approx0.71199$。该值明显低于正常区间 $[0.8,1.2]$，说明整车悬浮系统存在\textbf{功率放大器衰减型故障}，实际电磁力相对于理想电磁力整体衰减约 $28.80\%$。平滑窗口、边界样本剔除和首尾端点剔除等检验均表明该结论具有较好的稳健性。

针对问题四，在功率放大系数时变且相互独立条件下，单时刻仅有一个总电磁力观测方程，无法唯一反演16台电磁铁瞬时功率放大系数。因此，本文在短时间窗口内假设其近似不变，建立\textbf{滑动窗口正则回归模型}，估计等效功率放大系数，并结合\textbf{相对特征法}验证。结果表明，附件4整车总电磁力层面\textbf{无整体性功率放大故障}，但存在多个单电磁铁局部故障：$3\,\mathrm{s}\sim5\,\mathrm{s}$ 内，$1,3,14$ 号为衰减型故障，$7,15,16$ 号为突增型故障；$7\,\mathrm{s}\sim9\,\mathrm{s}$ 内，$2,3$ 号为衰减型故障，$16$ 号为突增型故障。综上，本文模型能够实现\textbf{电磁力拟合、动力学反演与故障检测定位}的一体化分析，可为磁浮列车悬浮系统状态监测提供参考。


 

